% PDF page 14
\source{gilbarg-trudinger:page-0014}
\chapter{Introduction}

\section*{Summary}

The principal objective of this work is the systematic development of the general
theory of second order quasilinear elliptic equations and of the linear theory
required in the process. This means we shall be concerned with the solvability of
boundary value problems (primarily the Dirichlet problem) and related general
properties of solutions of linear equations

\begin{equation}\label{gt.eq1.1}\dcref{(1.1)}\group{chapter-1}\source{gilbarg-trudinger:page-0014}
(1.1)\qquad Lu \equiv a^{ij}(x)D_{ij}u+b^i(x)D_i u+c(x)u=f(x),\qquad i,j=1,2,\dots,n,
\end{equation}

and of quasilinear equations

\begin{equation}\label{gt.eq1.2}\dcref{(1.2)}\group{chapter-1}\source{gilbarg-trudinger:page-0014}
(1.2)\qquad Qu \equiv a^{ij}(x,u,Du)D_{ij}u+b(x,u,Du)=0.
\end{equation}

Here $Du=(D_1u,\dots,D_nu)$, where $D_i u = \partial u/\partial x_i$, $D_{ij}u = \partial^2u/\partial x_i\,\partial x_j$, etc., and the
summation convention is understood. The ellipticity of these equations is expressed
by the fact that the coefficient matrix $[a^{ij}]$ is (in each case) positive definite in the
domain of the respective arguments. We refer to an equation as \emph{uniformly elliptic}
if the ratio $\gamma$ of maximum to minimum eigenvalue of the matrix $[a^{ij}]$ is bounded.
We shall be concerned with both non-uniformly and uniformly elliptic equations.
The classical prototypes of linear elliptic equations are of course Laplace's
equation

\begin{equation}\label{gt.eq1.3}\dcref{(1.3)}\group{chapter-1}\source{gilbarg-trudinger:page-0014}
\Delta u = \sum D_{ii}u = 0
\end{equation}

and its inhomogeneous counterpart, Poisson's equation $\Delta u=f$. Probably the best
known example of a quasilinear elliptic equation is the minimal surface equation

\begin{equation}\label{gt.eq1.4}\dcref{(1.4)}\group{chapter-1}\source{gilbarg-trudinger:page-0014}
\sum D_i\!\left(D_i u/(1+|Du|^2)^{1/2}\right)=0,
\end{equation}

which arises in the problem of least area. This equation is non-uniformly elliptic,
with $\gamma=1+|Du|^2$. The properties of the differential operators in these examples
motivate much of the theory of the general classes of equations discussed in this
book.

% PDF page 15
\source{gilbarg-trudinger:page-0015}

The relevant linear theory is developed in Chapters 2--9 (and in part of Chapter
12). Although this material has independent interest, the emphasis here is on
aspects needed for application to nonlinear problems. Thus the theory stresses weak
hypotheses on the coefficients and passes over many of the important classical and
modern results on linear elliptic equations.

Since we are ultimately interested in classical solutions of equation (1.2), what is
required at some point is an underlying theory of classical solutions for a suffi-
ciently large class of linear equations. This is provided by the Schauder theory in
Chapter 6, which is an essentially complete theory for the class of equations (1.1)
with Hölder continuous coefficients. Whereas such equations enjoy a definitive
existence and regularity theory for classical solutions, corresponding results cease
to be valid for equations in which the coefficients are assumed only continuous.

A natural starting point for the study of classical solutions is the theory of
Laplace's and Poisson's equations. This is the content of Chapters 2 and 4. In
anticipation of later developments the Dirichlet problem for harmonic functions
with continuous boundary values is approached through the Perron method of
subharmonic functions. This emphasizes the maximum principle, and with it the
barrier concept for studying boundary behavior, in arguments that are readily
extended to more general situations in later chapters. In Chapter 4 we derive the
basic Hölder estimates for Poisson's equation from an analysis of the Newtonian
potential. The principal result here (see Theorems 4.6, 4.8) states that all $C^2(\Omega)$
solutions of Poisson's equation, $\Delta u=f$, in a domain $\Omega$ of $\mathbb{R}^n$ satisfy a uniform
estimate in any subset $\Omega' \Subset \Omega$

\begin{equation}\label{gt.eq1.5}\dcref{(1.3)}\group{chapter-1}\source{gilbarg-trudinger:page-0015}
(1.3)\qquad \|u\|_{C^{2,\alpha}(\bar{\Omega}')} \leq C\!\left(\sup_{\Omega}|u|+\|f\|_{C^{\alpha}(\bar{\Omega})}\right),
\end{equation}

where $C$ is a constant depending only on $\alpha$ ($0<\alpha<1$), the dimension $n$ and dist
$(\Omega',\partial\Omega)$; (for notation see Section 4.1). This \emph{interior estimate} (interior since
$\Omega' \Subset \Omega$) can be extended to a \emph{global estimate} for solutions with sufficiently
smooth boundary values provided the boundary $\partial\Omega$ is also sufficiently smooth.
In Chapter 4 estimates up to the boundary are established only for hyperplane and
spherical boundaries, but these suffice for the later applications.

The climax of the theory of classical solutions of linear second order elliptic
equations is achieved in the Schauder theory, which is developed in modified and
expanded form in Chapter 6. Essentially, this theory extends the results of potential
theory to the class of equations (1.1) having Hölder continuous coefficients. This is
accomplished by the simple but fundamental device of regarding the equation
locally as a perturbation of the constant coefficient equation obtained by fixing the
leading coefficients at their values at a single point. A careful calculation based on
the above mentioned estimates for Poisson's equation yields the same inequality
(1.3) for any $C^{2,\alpha}$ solution of (1.1), where the constant $C$ now depends also on the
bounds and Hölder constants of the coefficients and in addition on the minimum
and maximum eigenvalues of the coefficient matrix $[a^{ij}]$ in $\Omega$. These results are
stated as interior estimates in terms of weighted interior norms (Theorem 6.2) and,
in the case of sufficiently smooth boundary data, as global estimates in terms of

% PDF page 16
\source{gilbarg-trudinger:page-0016}

global norms (Theorem 6.6). Here we meet the important and recurring concept of
an apriori estimate; namely, an estimate (in terms of given data) valid for all
possible solutions of a class of problems even if the hypotheses do not guarantee
the existence of such solutions. A major part of this book is devoted to the estab-
lishment of apriori bounds for various problems. (We have taken the liberty of
replacing the latin $a priori$ with the single word apriori, which will be used
throughout.)

    The importance of such apriori estimates is visible in several applications in
Chapter 6, among them in establishing the solvability of the Dirichlet problem by
the method of continuity (Theorem 6.8) and in proving the higher order regularity
of $C^2$ solutions under appropriate smoothness hypotheses (Theorems 6.17, 6.19).
In both cases the estimates provide the necessary compactness properties for
certain classes of solutions, from which the desired results are easily inferred.

    We remark on several additional features of Chapter 6, which are not needed
for the later developments but which broaden the scope of the basic Schauder
theory. In Section 6.5 it is seen that for continuous boundary values and a suitably
wide class of domains the proof of solvability of the Dirichlet problem for (1.1) can
be achieved entirely with interior estimates, thereby simplifying the structure of
the theory. The results of Section 6.6 extend the existence theory for the Dirichlet
problem to certain classes of non-uniformly elliptic equations. Here we see how
relations between geometric properties of the boundary and the degenerate ellipti-
city at the boundary determine the continuous assumption of boundary values.
The methods are based on barrier arguments that foreshadow analogous (but
deeper) results for nonlinear equations in Part II. In Section 6.7 we extend the
theory of (1.1) to the regular oblique derivative problem. The method is basically
an extrapolation to these boundary conditions of the earlier treatment of Poisson's
equation and the Schauder theory (without barrier arguments, however).

    In the preceding considerations, especially in the existence theory and barrier
arguments, the maximum principle for the operator $L$ (when $c \leq 0$) plays an
essential part. This is a special feature of second order elliptic equations that
simplifies and strengthens the theory. The basic facts concerning the maximum
principle, as well as illustrative applications of comparison methods, are contained
in Chapter 3. The maximum principle provides the earliest and simplest apriori
estimates of the general theory. It is of considerable interest that all the estimates of
Chapters 4 and 6 can be derived entirely from comparison arguments based on the
maximum principle, without any mention of the Newtonian potential or integrals.

    An alternative and more general approach to linear problems, without poten-
tial theory, can be achieved by Hilbert space methods based on generalized or
weak solutions, as in Chapter 8. To be more specific, let $L'$ be a second order
differential operator, with principal part of divergence form, defined by

\begin{equation}\label{gt.eq1.6}\dcref{(1.6)}\group{chapter-1}\source{gilbarg-trudinger:page-0017}
L'u \equiv D_i(a^{ij}(x)D_j u + b^i(x)u) + c^i(x)D_i u + d(x)u.
\end{equation}

If the coefficients are sufficiently smooth, then clearly this operator falls within the
class discussed in Chapter 6. However, even if the coefficients are in a much wider

% PDF page 17
\source{gilbarg-trudinger:page-0017}

class and $u$ is only weakly differentiable (in the sense of Chapter 7), one can still define weak or generalized solutions of $L'u=g$ in appropriate function classes. In particular, if the coefficients $a^{ij}$, $b^i$, $c^i$ are bounded and measurable in $\Omega$ and $g$ is an integrable function in $\Omega$, let us call $u$ a weak or generalized solution of $L'u=g$ in $\Omega$ if $u\in W^{1,2}(\Omega)$ (as defined in Chapter 7) and

\begin{equation}\label{gt.eq1.7}\dcref{(1.4)}\group{chapter-1}\source{gilbarg-trudinger:page-0017}
\int_{\Omega} \bigl[(a^{ij}D_j u+b^i u)D_i v-(c^i D_i u+d u)v\bigr]\,dx
=-\int_{\Omega} g v\,dx
\end{equation}

for all test functions $v\in C_0^1(\Omega)$. It is clear that if the coefficients and $g$ are sufficiently smooth and $u\in C^2(\Omega)$, then u is also a classical solution.

We can now speak also of weak solutions $u$ of the generalized Dirichlet problem,

\[
L'u=g \text{ in } \Omega,\qquad u=\varphi \text{ on } \partial\Omega,
\]

if $u$ is a weak solution satisfying $u-\varphi\in W_0^{1,2}(\Omega)$, where $\varphi\in W^{1,2}(\Omega)$. Assuming that the minimum eigenvalue of $[a^{ij}]$ is bounded away from zero in $\Omega$, that

\begin{equation}\label{gt.eq1.8}\dcref{(1.5)}\group{chapter-1}\source{gilbarg-trudinger:page-0017}
D_i b^i+d\leq 0
\end{equation}

in the weak sense, and that also $g\in L^2(\Omega)$, we find in Theorem 8.3 that the generalized Dirichlet problem has a unique solution $u\in W^{1,2}(\Omega)$. Condition (1.5), which is the analogue of $c\leq 0$ in (1.1), assures a maximum principle for weak solutions of $L'u\geq 0(\leq 0)$ (Theorem 8.1) and hence uniqueness for the generalized Dirichlet problem. Existence of a solution then follows from the Fredholm alternative for the operator $L'$ (Theorem 8.6), which is proved by an application of the Riesz representation theorem in the Hilbert space $W_0^{1,2}(\Omega)$.

The major part of Chapter 8 is taken up with the regularity theory for weak solutions. Additional regularity of the coefficients in (1.4) implies that the solutions belong to higher $W^{k,2}$ spaces (Theorems 8.8, 8.10). It follows from the Sobolev imbedding theorems in Chapter 7 that weak solutions are in fact classical solutions provided the coefficients are sufficiently regular. Global regularity of these solutions is inferred by extending interior regularity to the boundary when the boundary data are sufficiently smooth (Theorems 8.13, 8.14).

The regularity theory of weak solutions and the associated pointwise estimates are fundamental to the nonlinear theory. These results provide the starting point for the “bootstrap” arguments that are typical of nonlinear problems. Briefly, the idea here is to start with weak solutions of a quasilinear equation, regarding them as weak solutions of related linear equations obtained by inserting them into the coefficients, and then to proceed by establishing improved regularity of these solutions. Starting anew with the latter solutions and repeating the process, still further regularity is assured, and so on, until the original weak solutions are finally proved to be suitably smooth. This is the essence of the regularity proofs for the older variational problems and is implicit in the nonlinear theory presented here.

The Hölder estimates for weak solutions that are so vital for the nonlinear theory are derived in Chapter 8 from Harnack inequalities based on the Moser iteration technique (Theorems 8.17, 8.18, 8.20, 8.24). These results generalize the

% PDF page 18
\source{gilbarg-trudinger:page-0018}


basic a priori Hölder estimate of De Giorgi, which provided the initial breakthrough
in the theory of quasilinear equations in more than two independent variables.
The arguments rest on integral estimates for weak solutions $u$ derived from judi-
cious choice of test functions $v$ in (1.4). The test function technique is the dominant
theme in the derivation of estimates throughout most of this work.

In this edition we have added new material to Chapter 8 covering the Wiener
criterion for regular boundary points, eigenvalues and eigenfunctions, and Hölder
estimates for first derivatives of solutions of linear divergence structure equations.

We conclude Part I of the present edition with a new chapter, Chapter 9,
concerning strong solutions of linear elliptic equations. These are solutions which
possess second derivatives, at least in a weak sense, and satisfy (1.1) almost every-
where. Two strands are interwoven in this chapter. First we derive a maximum
principle of Aleksandrov, and a related a priori bound (Theorem 9.1) for solutions
in the Sobolev space $W^{2,n}(\Omega)$, thereby extending certain basic results from Chapter 3
to nonclassical solutions. Later in the chapter, these results are applied to establish
various pointwise estimates, including the recent Hölder and Harnack estimates of
Krylov and Safonov (Theorems 9.20, 9.22; Corollaries 9.24, 9.25). The other strand
in this chapter is the $L^p$ theory of linear second-order elliptic equations that is
analogous to the Schauder theory of Chapter 6. The basic estimate for Poisson’s
equation, namely the Calderon-Zygmund inequality (Theorem 9.9) is derived
through the Marcinkiewicz interpolation theorem, although without the use of
Fourier transform methods. Interior and global estimates in the Sobolev spaces
$W^{2,p}(\Omega)$, $1<p<\infty$, are established in Theorems 9.11, 9.13 and applied to the
Dirichlet problem for strong solutions, in Theorem 9.15 and Corollary 9.18.

Part II of this book is devoted largely to the Dirichlet problem and related
estimates for quasilinear equations. The results concern in part the general operator
(1.2) while others apply especially to operators of divergence form

\begin{equation}\label{gt.eq1.9}\dcref{(1.6)}\group{chapter-1}\source{gilbarg-trudinger:page-0018}
(1.6)\qquad Qu \equiv \diver A(x, u, Du) + B(x, u, Du)
\end{equation}

where $A(x, z, p)$ and $B(x, z, p)$ are respectively vector and scalar functions defined
on $\Omega \times \R \times \R^n$.

Chapter 10 extends maximum and comparison principles (analogous to results
in Chapter 3) to solutions and subsolutions of quasilinear equations. We mention
in particular a priori bounds for solutions of $Qu\geq 0\ (=0)$, where $Q$ is a divergence
form operator satisfying certain structure conditions more general than ellipticity
(Theorem 10.9).

Chapter 11 provides the basic framework for the solution of the Dirichlet
problem in the following chapters. We are concerned principally with classical
solutions, and the equations may be uniformly or non-uniformly elliptic. Under
suitable general hypotheses any globally smooth solution $u$ of the boundary value
problem for $Qu=0$ in a domain $\Omega$ with smooth boundary can be viewed as a
fixed point, $u = Tu$, of a compact operator $T$ from $C^{1,\alpha}(\overline{\Omega})$ to $C^{1,\alpha}(\overline{\Omega})$ for any
$\alpha \in (0,1)$. In the applications the function defined by $Tu$, for any $u \in C^{1,\alpha}(\overline{\Omega})$, is the
unique solution of the linear problem obtained by inserting $u$ into the coefficients
of $Q$. The Leray-Schauder fixed point theorem (proved in Chapter 11) then implies

% PDF page 19
\source{gilbarg-trudinger:page-0019}

    the existence of a solution of the boundary value problem provided an apriori
bound, in $C^{1,\alpha}(\overline{\Omega})$, can be established for the solutions of a related continuous
family of equations $u = T(u;\sigma)$, $0\leq \sigma \leq 1$, where $T(u;1)=Tu$ (Theorems 11.4,
11.6). The establishment of such bounds for certain broad classes of Dirichlet
problems is the object of Chapters 13-15.
    The general procedure for obtaining the required apriori bound for possible
solutions $u$ is a four-step process involving successive estimation of $\sup_{\partial\Omega}|u|$,
$\sup_{\Omega}|Du|$, $\sup_{\Omega}|D^2u|$, and $\|u\|_{C^{1,\alpha}(\overline{\Omega})}$ for some $\alpha>0$. Each of these estimates pre-
supposes the preceding ones and the final bound on $\|u\|_{C^{1,\alpha}(\overline{\Omega})}$ completes the
existence proof based on the Leray-Schauder theorem.
    As already observed, bounds on $\sup_{\Omega}|u|$ are discussed in Chapter 10. In the later
chapters this bound is either assumed in the hypotheses or is implied by properties
of the equation.
    Equations in two variables (Chapter 12) occupy a special place in the theory.
This is due in part to the distinctive methods that have been developed for them
and also to the results, some of which have no counterpart for equations in more
than two variables. The method of quasiconformal mappings and arguments based
on divergence structure equations (cf. Chapter 11) are both applicable to equations
in two variables and yield relatively easily the desired $C^{1,\alpha}$ apriori estimates, from
which a solution of the Dirichlet problem follows readily.
    Of particular interest is the fact that solutions of uniformly elliptic linear equa-
tions in two variables satisfy an apriori $C^{1,\alpha}$ estimate depending only on the
ellipticity constants and bounds on the coefficients, without any regularity assump-
tions (Theorem 12.4). Such a $C^{1,\alpha}$ estimate, or even the existence of a gradient
bound under the same general conditions is unknown for equations in more than
two variables. Another special feature of the two-dimensional theory is the
existence of an apriori $C^1$ bound $|Du|\leq K$ for solutions of arbitrary elliptic equations

\begin{equation}\label{gt.eq1.10}\dcref{(1.7)}\group{chapter-1}\source{gilbarg-trudinger:page-0019}
au_{xx}+2bu_{xy}+cu_{yy}=0,
\end{equation}

where $u$ is continuous on the closure of a bounded convex domain $\Omega$ and has
boundary value $\varphi$ on $\partial\Omega$ satisfying a bounded slope (or three-point) condition
with constant $K$. This classical result, usually based on a theorem of Rado on
saddle surfaces, is given an elementary proof in Lemma 12.6. The stated gradient
bound, which is valid for all solutions $u$ of the general quasilinear equation (1.7)
in which $a=a(x,y,u,u_x,u_y)$, etc., and such that $u=\varphi$ on $\hat{\partial}\Omega$, reduces this Dirichlet
problem to the case of uniformly elliptic equations treated in Theorem 12.5. In
Theorem 12.7 we obtain a solution of the general Dirichlet problem for (1.7),
assuming local Hölder continuity of the coefficients and a bounded slope condition
for the boundary data (without further smoothness restrictions on the data).
    Chapters 13, 14 and 15 are devoted to the derivation of the gradient estimates
involved in the existence procedure described above. In Chapter 13, we prove the
fundamental results of Ladyzhenskaya and Ural'tseva on Hölder estimates of
derivatives of elliptic quasilinear equations. In Chapter 14 we study the estimation
of the gradient of solutions of elliptic quasilinear equations on the boundary.

% PDF page 20
\source{gilbarg-trudinger:page-0020}

Summary

After considering general and convex domains, we give an account of the theory of Serrin which associates generalized boundary curvature conditions with the solvability of the Dirichlet problem. In particular, we are able to conclude from the results of Chapters 11, 13 and 14 the Jenkins and Serrin criterion for solvability of the Dirichlet problem for the minimal surface equation, namely, that this problem is solvable for smooth domains and arbitrary smooth boundary values if and only if the mean curvature of the boundary (with respect to the inner normal) is nonnegative at every point (Theorem 14.14).

Global and interior gradient bounds for solutions $u$ of quasilinear equations are established in Chapter 15. Following a refinement of an old procedure of Bernstein we derive estimates for $\sup_{\Omega}|Du|$ in terms of $\sup_{\partial\Omega}|Du|$ for classes of equations that include both uniformly elliptic equations satisfying natural growth conditions and equations sharing common structural properties with the prescribed mean curvature equation (Theorem 15.2). A variant of our approach yields interior gradient estimates for a more restricted class of equations (Theorem 15.3). We also consider uniformly and non-uniformly elliptic equations in divergence form (Theorems 15.6, 15.7 and 15.8), in which cases, by appropriate test function arguments, we deduce gradient estimates under different types of coefficient conditions than in the general case. We conclude Chapter 15 with a selection of existence theorems, chosen to illustrate the scope of the theory. These theorems are all obtained by various combinations of the apriori estimates in Chapters 10, 14 and 15 and a judicious choice of a related family of problems to which Theorem 11.8 can be applied.

In Chapter 16, we concentrate on the prescribed mean curvature equation and derive an interior gradient bound (Theorem 16.5) thereby enabling us to deduce existence theorems for the Dirichlet problem when only continuous boundary values are assigned (Theorems 16.8, 16.10). We also consider a family of equations in two variables, which in a certain sense bear the same relationship to the prescribed mean curvature equation as the uniformly elliptic equations of Chapter 12 bear to Laplace's equation. Indeed, by means of a generalized notion of quasiconformal mapping, we derive interior estimates for first and second derivatives. The second derivative estimates provide a generalization of a well known curvature estimate of Heinz for solutions of the minimal surface equation (Theorem 16.20) and moreover, imply an extension of the famous result of Bernstein that entire solutions of the minimal surface equation in two variables must be linear (Corollary 16.19). However, perhaps the striking feature of Theorems 16.5 and 16.20 is the approach. Rather than working in the domain $\Omega$, we work on the hypersurface $S$ given by the graph of the solution $u$ and exploit various relations between the tangential gradient and Laplacian operators on $S$ and the mean curvature of $S$.

We have also added to the present edition a new final chapter. Chapter 17 is devoted to fully nonlinear elliptic equations, which incorporates recent work on equations of Monge-Ampère and Bellman-Pucci type. These are equations of the general form

\begin{equation}\label{gt.eq1.11}\dcref{(1.8)}\group{chapter-1}\source{gilbarg-trudinger:page-0020}
F[u] = F(x,u,Du,D^2u) = 0
\end{equation}

% PDF page 21
\source{gilbarg-trudinger:page-0021}

and include linear and quasilinear equations of the forms (1.1) and (1.2) as special cases. The function $F$ is defined for $(x,z,p) \in \Omega \times \R \times \R^n \times \R^{n \times n}$ where $\R^{n \times n}$ denotes the linear space of real symmetric $n \times n$ matrices. Equation (1.8) is elliptic when the derivative $F_r$ is a positive definite matrix. The method of continuity (Theorem 17.8) reduces the solvability of the Dirichlet problem for (1.8) to the establishment of $C^{2,\alpha}(\bar{\Omega})$ estimates, for some $\alpha > 0$; that is, in addition to the first derivative estimation required for the quasilinear case, we need second derivative estimates for fully nonlinear equations. Such estimates are established for equations in two variables (Theorems 17.9, 17.10), uniformly elliptic equations (Theorems 17.14, 17.15) and equations of Monge-Amp\`ere type (Theorems 17.19, 17.20, 17.26), yielding, in particular, recent results on the solvability of the Dirichlet problem for uniformly elliptic equations by Evans, Krylov and Lions (in Theorems 17.17, 17.18), and for equations of Monge-Amp\`ere type by Krylov, and Caffarelli, Nirenberg and Spruck (in Theorem 17.23).

We conclude this summary with some guides to the reader. The material is not in strict logical order. Thus the theory of Poisson's equation (Chapter 4) would normally follow Laplace's equation (Chapter 2). However, the elementary character of the results on the maximum principle (Chapter 3) and the opportunity for the reader to meet early some general problems with variable coefficients recommends its insertion after Chapter 2. In fact, the general maximum principle is not used until the existence theory of Chapter 6. The basic material on functional analysis (Chapter 5) is needed in only a minor way for the Schauder theory: the contraction mapping principle and the basic concepts of Banach spaces suffice, except for the proof of the alternative in Theorem 6.15. For applications to nonlinear problems in Part II it is sufficient to know the results of Section 1--3 of Chapter 6. Depending on the reader's interests, it may be preferable to study the linear theory by starting directly with $L^2$ theory in Chapter 8; this assumes the preliminary material on functional analysis (Chapter 5) and on the calculus of weakly differentiable functions (Chapter 7). The Harnack inequalities and H\"older estimates in the regularity theory of Chapter 8 are not applied until Chapter 13.

    The theory of quasilinear equations in two variables (Chapter 12) is essentially independent of Chapters 7--11 and can be read following Chapter 6 provided one assumes the Schauder fixed point theorem (Theorem 11.1). The method of quasi-conformal mappings is met again in Chapter 16 but otherwise the remaining chapters are independent of Chapter 12. Accordingly, after the basic outline of the nonlinear theory in Chapter 11 the reader can proceed directly to the $n$-variable theory in Chapters 13--17. Chapter 16 is largely independent of Chapters 13--15. Chapters 6 and 9 are sufficient preparation for most of Chapter 17.

\textbf{Further Remarks}

Beyond the assumption of basic real analysis and linear algebra the material in this work is almost entirely self-contained. Thus, much of the preliminary development


% PDF page 22
\source{gilbarg-trudinger:page-0022}



of potential theory and functional analysis, as well as results on Sobolev spaces and
fixed point theorems, will be familiar to many readers, although the proof of the
Leray-Schauder theorem without topological degree in Theorem 11.6 is probably
not so well known. A number of well established auxiliary results, such as the
interpolation inequalities and extension lemmas of Chapter 6, are proved for the
sake of completeness.

There is substantial overlap with the monographs of Ladyzhenskaya and
Ural'tseva [LU 4] and Morrey [MY 5]. This book differs from the former in some
of the analytical techniques and in the emphasis on non-uniformly elliptic equations
in the nonlinear theory; it differs from the latter in not being directly concerned
with variational problems and methods. The present work also includes material
developed since the publication of those books. On the other hand, it is much more
limited in various ways. Among the topics not included are systems of equations,
semilinear equations, the theory of monotone operators, and aspects of the subject
based on geometric measure theory.

In a subject that is often quite technical we have not always striven for the
greatest generality, especially with respect to the modulus of continuity, estimates,
integral conditions, and the like. We have instead confined ourselves to conditions
determined by \emph{power} functions: for example, H\"older continuity rather than Dini
continuity, $L^p$ spaces in Chapter 8 rather than Orlicz spaces, structure conditions
in terms of powers of $|p|$ rather than more general functions of $|p|$, etc. By suitable
modification of the proofs the reader will usually be able to supply the appropriate
generalizations.

Historical material and bibliographical references are discussed primarily in
the Notes at the end of the chapters. These are not intended to be complete but
rather to supplement the text and place it in better perspective. A much more
extensive survey of the literature until 1968 is contained in Miranda [MR 2]. The
problems attached to the chapters are also intended to supplement the text;
hopefully they will be useful exercises for the reader.

\bigskip

\begin{notation}\label{gt.not.basic-notation}\dcref{Basic Notation}\group{chapter-1}\source{gilbarg-trudinger:page-0022,gilbarg-trudinger:page-0023}
\noindent\textbf{Basic Notation}

\begin{description}
\item[$\R^n$:] Euclidean $n$-space, $n \geq 2$, with points $x = (x_1,\dots,x_n)$, $x_i \in \R$ (real numbers);
$|x| = \left(\sum x_i^2\right)^{1/2}$; if $b = (b_1,\dots,b_n)$ is an ordered $n$-tuple, then
$\|b\| = \left(\sum b_i^2\right)^{1/2}$.
\item[$\R^n_+$:] half-space in $\R^n$ $= \{x \in \R^n \mid x_n > 0\}$.
\item[$\partial S$:] boundary of the point set $S$; $\overline{S} = \closure{S} = S \cup \partial S$.
\item[$S - S'$:] $\{x \in S \mid x \notin S'\}$.
\item[$S' \subset\subset S$:] $S'$ has compact closure in $S$; $S'$ is strictly contained in $S$.
\item[$\Omega$:] a proper open subset of $\R^n$, not necessarily bounded; $\Omega$ is a domain if it is also connected; $|\Omega|$ = volume of $\Omega$.
\item[$B(y)$:] a ball in $\R^n$ with center $y$; $B_r(y)$ is the open ball of radius $r$ centered at $y$.
\item[$\omega_n$:] volume of unit ball in $\R^n$ $= \dfrac{2\pi^{n/2}}{n\Gamma(n/2)}$.
\item[$D_i u = \partial u/\partial x_i$, $D_{ij}u = \partial^2 u/\partial x_i \partial x_j$, etc.:] $Du = (D_1u,\dots,D_nu) = \mathrm{gradient\ of\ }u$;
$D^2u = [D_{ij}u]$ = Hessian matrix of second derivatives $D_{ij}u$, $i,j = 1,2,\dots,n$.
\end{description}
\end{notation}


% PDF page 23
\source{gilbarg-trudinger:page-0023}

\[
\beta=(\beta_1,\dots,\beta_n),\ \beta_i=\text{integer}\ge 0,\ \text{with }|\beta|=\sum \beta_i,\ \text{is a multi-index; we define}
\]
\[
D^\beta u=\frac{\partial^{|\beta|}u}{\partial x_1^{\beta_1}\cdots \partial x_n^{\beta_n}}
\]

\[
C^0(\Omega)\ \bigl(C^0(\overline{\Omega})\bigr): \text{the set of continuous functions on }\Omega\ (\overline{\Omega}).
\]
\[
C^k(\Omega): \text{the set of functions having all derivatives of order }\leq k\text{ continuous in }
\]
\[
\Omega\ (k=\text{integer}\ge 0\text{ or }k=\infty).
\]
\[
C^k(\overline{\Omega}): \text{the set of functions in }C^k(\Omega)\text{ all of whose derivatives of order }\leq k\text{ have}
\]
\[
\text{continuous extensions to }\overline{\Omega}.
\]
\[
\supp\,u: \text{the support of }u,\text{ the closure of the set on which }u\neq 0.
\]
\[
C_0^k(\Omega): \text{the set of functions in }C^k(\Omega)\text{ with compact support in }\Omega.
\]
\[
C=C(*,\dots,*)\text{ denotes a constant depending only on the quantities appearing}
\]
\[
\text{in parentheses. In a given context, the same letter }C\text{ will (generally) be used to}
\]
\[
\text{denote different constants depending on the same set of arguments.}
\]

% PDF page 24
\source{gilbarg-trudinger:page-0024}

Part I

Linear Equations

% PDF page 25
\source{gilbarg-trudinger:page-0025}

\begin{center}
\textbf{Chapter 2}

\vspace{1.2em}

\textbf{\LARGE Laplace's Equation}
\end{center}

Let $\Omega$ be a domain in $\R^n$ and $u$ a $C^2(\Omega)$ function. The Laplacian of $u$, denoted $\Delta u$, is defined by
\begin{equation}
\tag{2.1}
\Delta u=\sum_{i=1}^n D_{ii}u=\diver Du.
\end{equation}

The function $u$ is called \emph{harmonic} (\emph{subharmonic}, \emph{superharmonic}) in $\Omega$ if it satisfies there
\begin{equation}
\tag{2.2}
\Delta u=0 \qquad (\geq 0, \leq 0).
\end{equation}

In this chapter we develop some basic properties of harmonic, subharmonic and superharmonic functions which we use to study the solvability of the classical Dirichlet problem for \emph{Laplace's equation}, $\Delta u=0$. As mentioned in Chapter 1, Laplace's equation and its inhomogeneous form, Poisson's equation, are basic models of linear elliptic equations.

Our starting point here will be the well known divergence theorem in $\R^n$. Let $\Omega_0$ be a bounded domain with $C^1$ boundary $\partial\Omega_0$ and let $\nu$ denote the unit outward normal to $\partial\Omega_0$. For any vector field $\mathbf{w}$ in $C^1(\overline{\Omega}_0)$, we then have
\begin{equation}
\tag{2.3}
\int_{\Omega_0} \diver \mathbf{w}\, dx=\int_{\partial\Omega_0} \mathbf{w}\cdot \nu\, ds
\end{equation}
where $ds$ indicates the $(n-1)$-dimensional area element in $\partial\Omega_0$. In particular if $u$ is a $C^2(\overline{\Omega}_0)$ function we have, by taking $\mathbf{w}=Du$ in (2.3),
\begin{equation}
\tag{2.4}
\int_{\Omega_0} \Delta u\, dx=\int_{\partial\Omega_0} Du\cdot \nu\, ds=\int_{\partial\Omega_0} \frac{\partial u}{\partial \nu}\, ds.
\end{equation}

(For a more general formulation of the divergence theorem, see [KE 2].)

\vspace{1.5em}

\noindent\textbf{2.1. The Mean Value Inequalities}

Our first theorem, which is a consequence of the identity (2.4), comprises the well known mean value properties of harmonic, subharmonic and superharmonic functions.
